% =====================================================================
%  Trecho para a carta-resposta — pgsg_1 (CHEMOLAB-D-26-00658)
%  Correções identificadas pelos próprios autores após a submissão da R1.
%  Inserir como seção inicial, ANTES das respostas ponto a ponto.
% =====================================================================

\section*{Author-initiated corrections}

While preparing the extension of this work to a second spectroscopic
modality, we re-examined the experimental pipeline of the present manuscript
and identified a replication error that affects several reported quantities.
We report it here in full, together with the corrected results, before
addressing the reviewers' individual comments. The correction does not alter
the central claim of the paper, but it changes the magnitude of several
secondary results, retracts one claim made in the abstract, and --- in
tracing the error to its consequences --- produced a new and stronger result
in its place (Section~3 below).

\subsection*{1. Degenerate replicates at the largest training size}

In the experiment runner, the model factories fixed the random seed at 42:

\begin{verbatim}
lambda: PGSGv2Model(hidden=32, max_epochs=500, patience=30, seed=42)
\end{verbatim}

The five ``seeds'' therefore varied only the draw of the training subset
performed by the stratified sampler, never the initialisation of the model
itself. At $n = 1{,}159$ the requested training size coincides with the size
of the training pool, so the sampler returns the same subset for every seed;
combined with a test split fixed at seed 42 and a model seed fixed at 42,
the five replicates at that point are identical runs. We verified this
directly in the saved output: at $n_{\text{actual}} = 1{,}158$ the five
recorded $R^2$ values agree to the sixth decimal place for all five models
(PGSGv2 $= 0.805524$, PGSGv2-random $= 0.757743$, PLS $= 0.568351$). At
smaller training sizes, where the requested size is below the pool size, the
values do vary across seeds, which confirms the mechanism.

Two consequences follow. First, the bootstrap confidence intervals reported
for the $n = 1{,}159$ row of Table~1 were computed by resampling five copies
of a single measurement, which is why they appear as zero-width intervals
(e.g.\ $[0.81, 0.81]$ and $[0.76, 0.76]$). They are not confidence intervals
and have been removed. Second, and more consequentially, the reported
$\Delta_{\text{prior}} = +0.048$ at that point is a single measurement
against another single measurement, not a mean over five replicates.

We have re-executed the full scale curve with ten seeds, each varying both
the training-subset draw and the model initialisation, at every point of the
grid. The revised Tables~1 and~2 report these results. As an internal check,
PLS---being deterministic---remains identical across seeds at
$n = 1{,}158$, as expected, and reproduces the previously published value
exactly.

\subsection*{2. Revised result: the literature prior does not improve accuracy}

Under genuine replication, the paired difference between the
literature-initialised and uninformed conditions is not significant at any
training size, and its sign is not consistent:

\begin{center}
\small
\begin{tabular}{rcc}
\hline
$n$ & $\Delta_{\text{prior}}$ (paired) & $p$ (Wilcoxon) \\
\hline
30      & $+0.010 \pm 0.078$ & $0.625$ \\
60      & $-0.013 \pm 0.045$ & $0.322$ \\
100     & $-0.046 \pm 0.060$ & $0.037$ \\
200     & $+0.027 \pm 0.051$ & $0.193$ \\
400     & $+0.011 \pm 0.051$ & $0.695$ \\
700     & $+0.019 \pm 0.028$ & $0.084$ \\
1{,}158 & $-0.005 \pm 0.013$ & $0.770$ \\
\hline
\end{tabular}
\end{center}

The only nominally significant value, at $n = 100$, is in the direction
\emph{opposite} to the prior, and is of the order expected by chance across
seven tests; we do not interpret it as evidence against the prior either.

We have accordingly revised the claims that the prior provides a measurable
advantage in data-scarce regimes and that it continues to act as a beneficial
regulariser at the largest training size. The correct statement is that the
literature prior yields predictive performance statistically indistinguishable
from an uninformed initialisation across the entire sample-size range studied.

We note that this revised conclusion agrees with an independent evaluation on
the same dataset reported in a companion study of cross-modality
generalisation, in which ten-seed ablations in both NIR and Raman likewise
found statistical equivalence in accuracy between the two initialisations,
while gate topology was markedly more reproducible under the literature prior
(pairwise Jaccard $0.90$--$0.94$ versus $0.58$--$0.65$). Taken together, the
two studies support a more precise characterisation of what the prior
contributes: interpretive stability rather than predictive gain. We consider
this a clarification of the contribution rather than a weakening of it---a
gate that selects different bands on every run is not a trustworthy
explanation, even when its predictions are accurate.

\subsection*{3. New result: accuracy and interpretability dissociate}

Having established that the prior does not improve accuracy, we re-measured
the two interpretability metrics --- which the previous version reported over
five seeds --- under ten seeds and, for the first time, under \emph{both}
initialisations. The gate vectors were extracted at every point of the grid
and are deposited with the revision.

The literature-prior figures reproduce those previously reported:
$\rho(\mathbf{g},\mathbf{s}) \geq 0.9987$ and pairwise Jaccard
$0.945$--$0.987$ across the sample-size range. The uninformed condition,
never previously measured on these metrics, differs sharply:

\begin{center}
\small
\begin{tabular}{r cc cc}
\hline
& \multicolumn{2}{c}{Literature prior} & \multicolumn{2}{c}{Uninformed} \\
$n$ & $\rho$ & Jaccard & $\rho$ & Jaccard \\
\hline
     30 & $1.0000$ & $0.960$ & $0.035$ & $0.112$ \\
     60 & $1.0000$ & $0.956$ & $0.025$ & $0.090$ \\
    100 & $1.0000$ & $0.969$ & $0.068$ & $0.166$ \\
    200 & $0.9998$ & $0.987$ & $0.192$ & $0.307$ \\
    400 & $0.9996$ & $0.976$ & $0.190$ & $0.433$ \\
    700 & $0.9995$ & $0.945$ & $0.238$ & $0.386$ \\
1{,}159 & $0.9987$ & $0.950$ & $0.280$ & $0.566$ \\
\hline
\end{tabular}
\end{center}

Two conditions that are statistically indistinguishable in predictive
accuracy at every sample size (Section~2 above) thus differ by a factor of
between $1.7$ and $10.6$ in the reproducibility of the band set they select,
and the uninformed gate never exceeds $\rho = 0.28$ against the literature
profile. The optimisation admits many solutions of equivalent accuracy whose
spectral weightings bear little relation to one another or to known
absorption chemistry; the prior does not find a better one, it selects the
one that can be read.

We consider this the most substantive gain of the revision. The original
manuscript claimed a small accuracy benefit from the prior, which did not
survive replication; in its place we can now report a large, consistent and
directly measured benefit of a different kind, on the property that
motivates the method in the first place. Section~5.3 and Table~3 of the
revised manuscript are organised around this result, and Table~2 now reports
performance separately from interpretability so that the dissociation is
legible at a glance.

\subsection*{4. Retracted claim in the abstract}

The abstract stated that for $n \leq 60$, PGSGv2 is the only method achieving
positive $R^2$, and attributed this to prior-guided initialisation. Under
ten-seed replication this is not correct on either count: at $n = 60$ the
shallow CNN also attains a positive mean ($+0.013$), and the uninformed
variant PGSGv2-random is positive at both $n = 30$ ($0.073$) and $n = 60$
($0.138$), where it in fact slightly exceeds the literature-initialised
variant.

The underlying finding is nevertheless robust and remains the strongest
result of the paper: in the small-sample regime, PLS ($-0.53$ at $n = 30$,
$-0.08$ at $n = 60$) and the unconstrained MLP ($-3.05$ and $-1.41$) fail
severely, while both PGSGv2 variants remain positive with an order of
magnitude less variance. The advantage is attributable to the gated
architecture, not to the chemical content of the initialisation. We have
rewritten the sentence accordingly.

\subsection*{5. Revised baseline comparison and crossover point}

With genuine replication the ordering of the baselines at $n = 1{,}158$
changes: the MLP mean rises to $0.749$ and the shallow CNN falls to $0.690$,
making the MLP the strongest baseline. The relative gains of PGSGv2-lit are
therefore $+38\%$ over PLS, $+5\%$ over MLP and $+14\%$ over CNN1D, replacing
the previously reported $+42\%$, $+15\%$ and $+14\%$. These figures have been
updated in the abstract, Section~5.2 and the conclusion.

The crossover with PLS is also no longer a single point. PGSGv2-lit exceeds
PLS at $n = 30$ and $n = 60$, falls below it at $n = 100$ and $n = 200$, and
exceeds it consistently from $n = 400$ onward. The estimate $n^{*} \approx 95$
was obtained by interpolating between two points of a curve that is not
monotonic in this region. We have removed the interpolated crossover estimate
and replaced it with the direct statement that PGSGv2 is superior in the
small-sample regime and from $n \geq 400$, and comparable to or below PLS in
the intermediate range. Figure~1(b) has been redrawn accordingly.

\subsection*{6. Band-selection protocol}

Section~4.1 described the removal of all bands with zero variance in any of
the five seasons (the ``union-of-zeroed-bands'' protocol), yielding
$p = 281$. The code computes the mask from the training split of each
scenario instead. In the within-season configuration reported here the two
procedures coincide at $p = 281$ for every scenario, so no reported result is
affected; we verified this explicitly in the re-execution, where $p = 281$
holds at all seven training sizes and all ten seeds. We have nevertheless
corrected the description in Section~4.1 to state what the implementation
does, and we note in the text that a mask derived per split would not in
general be constant across the grid, which is a requirement for the scale
curve to be interpretable.

\subsection*{7. Typographical corrections}

Two citation commands lost their leading backslash and were typeset as
literal text in the compiled PDF (\texttt{citegoodfellow2016} in the
introduction and \texttt{citehoerl1970} in Section~3.3); both are corrected.
In Table~2 of the previous version, the $\Delta_{\text{PLS}}$ column carried
inverted signs relative to the text, to Figure~2(a) and to the values
computable from Table~1 (e.g.\ $-0.60$ tabulated at $n = 30$ where the text
and the definition in the caption give $+0.60$). The column has been
corrected.

\subsection*{Availability}

The re-execution scripts, the raw per-seed outputs, the extracted gate
vectors for every point of the grid, and the diagnostic that identified the
degenerate replicates are available in the public repository accompanying
this manuscript, alongside the original outputs, so that the correction can
be audited in full.
